% !TEX program = pdflatex
%
% citation-check-cost-audit.tex
%
% Full-fidelity transcript + analysis of the Citation Check cost audit
% conducted 2026-04-07 between Aditya Nittur (CEO/CoFounder, Flowblinq) and
% the CoFounder agent. Captures the file-level audit of
% geo/lib/services/citation-checker.ts, the provider pricing reverification,
% the architectural alternatives (external search, preemptive injection,
% prompt caching, MCP), and the final pricing decision.
%
% Output: geo/docs/citation-check-cost-audit.pdf
% Compile: pdflatex citation-check-cost-audit.tex  (run twice for ToC)

\documentclass[11pt,letterpaper]{article}

\usepackage[margin=0.9in]{geometry}
\usepackage{parskip}
\usepackage{booktabs}
\usepackage{longtable}
\usepackage{array}
\usepackage{xcolor}
\usepackage{colortbl}
\usepackage{hyperref}
\usepackage{fancyhdr}
\usepackage{enumitem}
\usepackage{listings}
\usepackage{framed}
\usepackage{titling}
\usepackage{sectsty}

\definecolor{brand}{HTML}{1F3A5F}
\definecolor{accent}{HTML}{B8733A}
\definecolor{codebg}{HTML}{F5F5F5}
\definecolor{codeborder}{HTML}{DDDDDD}
\definecolor{quoteblue}{HTML}{EEF2F8}
\definecolor{warn}{HTML}{C0392B}
\definecolor{good}{HTML}{1E8449}

\hypersetup{
  colorlinks=true,
  linkcolor=brand,
  urlcolor=brand,
  citecolor=brand,
  pdftitle={Citation Check Cost Audit},
  pdfauthor={CoFounder agent for Aditya Nittur},
}

\sectionfont{\color{brand}}
\subsectionfont{\color{brand}}
\subsubsectionfont{\color{brand}}

\pagestyle{fancy}
\fancyhf{}
\fancyhead[L]{\color{brand}\small Citation Check Cost Audit}
\fancyhead[R]{\color{brand}\small 2026-04-07}
\fancyfoot[C]{\thepage}
\renewcommand{\headrulewidth}{0.4pt}

\lstdefinestyle{code}{
  backgroundcolor=\color{codebg},
  basicstyle=\ttfamily\footnotesize,
  frame=single,
  rulecolor=\color{codeborder},
  breaklines=true,
  showstringspaces=false,
  captionpos=b,
}
\lstset{style=code}

\newenvironment{userturn}
  {\begin{quote}\small\color{brand}\textbf{Aditya:}~}
  {\end{quote}}

\newenvironment{cofounderturn}
  {\par\vspace{0.2em}\noindent\textbf{CoFounder:}\par}
  {\par\vspace{0.4em}}

\newcommand{\good}[1]{\textcolor{good}{#1}}
\newcommand{\warn}[1]{\textcolor{warn}{#1}}
\newcommand{\costneg}[1]{\textcolor{warn}{#1}}
\newcommand{\costpos}[1]{\textcolor{good}{#1}}

\title{\color{brand}\Huge Citation Check Cost Audit\\[0.4em]
  \Large \texttt{geo/lib/services/citation-checker.ts}}
\author{CoFounder agent\\for Aditya Nittur\\Flowblinq}
\date{2026-04-07}

\begin{document}

\maketitle

\begin{center}
\begin{minipage}{0.85\textwidth}
\small
\noindent\textbf{Abstract.} This document captures, in full fidelity, the
cost audit conversation conducted on 2026-04-07 between Aditya Nittur and the
CoFounder agent. The audit began as a deep-dive into the cost of running the
\texttt{citation-checker.ts} service in response to a user clicking the
``Run Citation Check'' button in the GEO product UI. It expanded into a
reverification of live LLM provider pricing, an analysis of architectural
alternatives (external search, preemptive tool injection, prompt caching,
Model Context Protocol), and culminated in a pricing decision at fixed base
subscriptions with metered overage priced at 1.5$\times$ variable cost.

\noindent The document is intended as (1) a decision record for the
architectural and pricing choices made, (2) the source material for the
follow-up technical specifications (TS-077 architecture, TS-078 Stripe
metered billing), and (3) evidence for investor or partner conversations
about Flowblinq's unit economics on the citation check feature.
\end{minipage}
\end{center}

\vspace{1em}

\tableofcontents

\newpage

%%%%%%%%%%%%%%%%%%%%%%%%%%%%%%%%%%%%%%%%%%%%%%%%%%%%%%%%%%%%%%%%%%%%%%%%%%%%%%%

\section{Context and opening prompt}

The session opened with the CoFounder agent picking up a crashed pipeline:
no agents were active except HolePoker (which had just restarted), and the
stale \texttt{ACTIVATE-1.yaml} signal was from a completed ES-075 run. The
last real activity was the TS-076 taskboarding on 2026-04-04.

Aditya's opening:

\begin{userturn}
Do a deep-dive audit of all costs when running citation-checker.ts --- all
of the costs.
\end{userturn}

The CoFounder clarified the interpretation: the file
\texttt{geo/lib/services/citation-checker.ts} is the service layer that runs
when a user clicks ``Run Citation Check'' in the product, and ``all of the
costs'' was interpreted to mean (a) direct LLM API spend, (b) third-party
search/tool fees, (c) infrastructure costs (Vercel, database, queue), and
(d) credit charges levied against the customer.

\section{What the file actually does}

\subsection{Cost-bearing surface}

\texttt{citation-checker.ts} exports one cost-bearing function
(\texttt{runCitationCheck}) and several pure aggregators
(\texttt{aggregateByDimension}, \texttt{aggregateCompetitorsByDimension},
\texttt{computeImpressionShare}, \texttt{generateTierInsight},
\texttt{generateDominanceInsights}). Only \texttt{runCitationCheck} makes
paid external calls.

Inside \texttt{runCitationCheck}, the entire external-cost surface is
\textbf{four helper functions}, one per provider:

\begin{center}
\begin{tabular}{@{}llllll@{}}
\toprule
\textbf{Helper} & \textbf{File line} & \textbf{Provider} & \textbf{Model} & \textbf{Tool} & \textbf{Output cap} \\
\midrule
\texttt{queryOpenAI}     & 85--110  & OpenAI    & gpt-5.4-mini               & web\_search          & 256 tok \\
\texttt{queryAnthropic}  & 112--130 & Anthropic & claude-haiku-4-5-20251001  & web\_search\_20250305 & 256 tok \\
\texttt{queryPerplexity} & 132--150 & Perplexity (via OpenAI shim) & sonar    & built-in retrieval   & 256 tok \\
\texttt{queryGoogle}     & 152--166 & Gemini    & gemini-2.5-flash           & googleSearch         & unset \\
\bottomrule
\end{tabular}
\end{center}

The file itself does \textbf{zero} database writes, \textbf{zero} Redis
operations, \textbf{zero} Blob storage writes, and \textbf{zero} queue or
CostMaster calls. Everything that costs money is concentrated in those four
functions.

\subsection{Call volume per invocation}

Volume = prompt count $\times$ provider count.

\textbf{Prompts} are generated upstream by
\texttt{citation-prompt-generator.ts}. The V2 path uses
\texttt{buildCoveringArray} with a budget of 36 indirect prompts, plus
\texttt{buildDirectPromptsV2} which returns 8 direct prompts. Total:
\textbf{44 prompts per run} in the typical V2 path, with a hard ceiling of
48 prompts under the legacy path.

\textbf{Providers} are returned by \texttt{getConfiguredProviders}
(lines 170--181). With all four environment keys set, there are 4
providers.

Cartesian product: $44 \times 4 = \mathbf{176}$ LLM calls per click of Run
Citation Check. Ceiling: $48 \times 4 = 192$.

Batching in \texttt{config.ts}: \texttt{CITATION\_CHECK\_BATCH\_SIZE = 20}
parallel calls, \texttt{CITATION\_CHECK\_BATCH\_DELAY\_MS = 100}. Nine batches
per run, wall time 50--150 seconds end-to-end.

\subsection{Per-call token and tool fees}

Each call passes approximately:

\begin{itemize}[nosep]
  \item System prompt \texttt{CITATION\_SYSTEM\_PROMPT} $\approx$ 250 input tokens (shared across providers)
  \item User prompt (single citation query) $\approx$ 20--40 tokens
  \item Tool schema (web\_search) $\approx$ 80--150 tokens
  \item \textbf{Search result injection when the tool fires: $+$1,500--4,000 input tokens per search --- the silent killer}
  \item Output capped at 256 tokens
\end{itemize}

Tool-usage caps:
\begin{itemize}[nosep]
  \item Anthropic: \texttt{max\_uses: 2} is a hard API-level cap
  \item OpenAI: \texttt{search\_context\_size: "low"} is a soft hint only; there is no hard cap on search count
  \item Gemini: grounding is fully provider-managed, no exposed cap
  \item Perplexity Sonar: search is baked into the model; no cap
\end{itemize}

\section{Initial cost rollup (conservative list-price math)}

\begin{center}
\begin{tabular}{@{}lrrrr@{}}
\toprule
\textbf{Provider} & \textbf{Calls} & \textbf{Mid \$/call} & \textbf{Mid total} & \textbf{High total} \\
\midrule
OpenAI               & 44 & \$0.027 & \$1.19 & \$1.50 \\
Anthropic            & 44 & \$0.026 & \$1.14 & \$1.40 \\
Perplexity           & 44 & \$0.007 & \$0.31 & \$0.40 \\
Gemini (under quota) & 44 & \$0.002 & \$0.09 & \$0.09 \\
Gemini (over quota)  & 44 & \$0.037 & \$1.63 & \$1.85 \\
\midrule
\textbf{Total (Gemini under quota)} & 176 & & \textbf{\$2.73} & \textbf{\$3.39} \\
\textbf{Total (Gemini over quota)}  & 176 & & \textbf{\$4.27} & \textbf{\$5.15} \\
\bottomrule
\end{tabular}
\end{center}

\subsection{What the customer pays (initial misreading)}

CoFounder initially quoted \$0.24/credit from \texttt{BULK\_CREDIT\_PRICE\_INR}
and reported customer payment as \$1.20/click (5 credits $\times$ \$0.24).
This was later corrected by Aditya. The bulk-INR price applies only to the
bulk CSV audit path, not to the per-action \texttt{shareOfVoice} charge,
which draws from the fungible subscription credit pool at much lower \$/credit
rates.

\subsection{Margin finding (before correction)}

Even with the incorrect \$1.20 revenue figure, the conclusion was:

\begin{quote}
\textbf{Citation Check is loss-making at the per-run unit level on plausible
assumptions.} The 5-credit price was set when the call surface was much
smaller; the V2 prompt architecture (44 prompts $\times$ 4 providers
$+$ web\_search) is materially more expensive than the credit model assumes.
\end{quote}

\section{Clarification: is an audit being called on this path?}

\begin{userturn}
Again, I am asking only about the Run Citation Check action from the UI,
are we calling an audit on that path?
\end{userturn}

CoFounder traced the UI button to \texttt{handleScanCitations} in
\texttt{SitePageClient.tsx:335} and \texttt{runCheck} in
\texttt{citation-monitor.tsx:493}, both of which \texttt{POST} to
\texttt{/api/sites/[id]/citation-check}. That route:

\begin{itemize}[nosep]
  \item Requires \texttt{site.geoScorecard} to already exist (line 45--49),
    returning \texttt{422 geo\_analysis\_required} if not --- so the route
    \textit{refuses to start} without an audit already present
  \item Deducts exactly \texttt{ACTION\_CREDITS.shareOfVoice = 5} credits
  \item Runs lazy one-shot enrichments (\texttt{extractTrees},
    \texttt{extractCategoriesViaHaiku}, \texttt{extractBrandKeywords}) that
    operate on already-crawled data and do not re-crawl or re-run the audit
\end{itemize}

\textbf{Conclusion:} clicking Run Citation Check does \textbf{not} trigger a
GEO audit. The two buttons are separate. The previous section's ``adjacent
cost'' calls all live inside the citation-check route and are part of the
citation check, not a separate audit double-charge.

\section{Reverification with live provider pricing}

\begin{userturn}
Check the current pricing and reverify. [Later clarified: ``I meant check
the GEO pricing and reverify the cost per citation that we are bearing.'']
\end{userturn}

CoFounder pulled live prices from vendor docs.

\subsection{GEO customer-facing pricing (from \texttt{geo/lib/config.ts})}

\begin{center}
\begin{tabular}{@{}llr@{}}
\toprule
\textbf{Source of credits} & \textbf{Math} & \textbf{\$/credit} \\
\midrule
One-off pack           & 100 credits / \$10     & \$0.1000 \\
Starter subscription   & 1,430 credits / \$99   & \$0.0692 \\
Growth subscription    & 7,465 credits / \$249  & \$0.0334 \\
Pro subscription       & 28,915 credits / \$499 & \$0.0173 \\
Bulk CSV (INR path)    & 1 credit / 20 INR $\approx$ \$0.24 & \$0.2400 \\
\bottomrule
\end{tabular}
\end{center}

Bulk INR pricing applies only to bulk CSV audits, not citation check.
Citation check's 5-credit charge yields customer revenue per click of:

\begin{center}
\begin{tabular}{@{}lr@{}}
\toprule
\textbf{Credit source} & \textbf{5 credits cost customer} \\
\midrule
Pro         & \$0.087 \\
Growth      & \$0.167 \\
Starter     & \$0.346 \\
One-off pack & \$0.500 \\
\bottomrule
\end{tabular}
\end{center}

\textbf{Prior ``\$1.20/click'' figure was 2.4$\times$ to 14$\times$ too high.}

\subsection{LLM unit prices (verified live)}

\begin{center}
\begin{tabular}{@{}llrr@{}}
\toprule
\textbf{Provider / Model} & \textbf{Tool} & \textbf{Input \$/M} & \textbf{Output \$/M} \\
\midrule
OpenAI \texttt{gpt-5.4-mini}         & web\_search \$10/1k + 8K search-content at input rate & \$0.75 & \$4.50 \\
Anthropic \texttt{claude-haiku-4-5}  & web\_search \$10/1k + search results at input rate    & \$1.00 & \$5.00 \\
Google \texttt{gemini-2.5-flash}     & Grounding free $\le$1,500/day, then \$35/1k            & \$0.30 & \$2.50 \\
Perplexity \texttt{sonar}            & \$5/1k requests at \texttt{low} context                & \$1.00 & \$1.00 \\
Anthropic \texttt{claude-sonnet-4-6} (ancillary) & ---                                        & \$3.00 & \$15.00 \\
\bottomrule
\end{tabular}
\end{center}

\subsection{Corrected per-call cost (current architecture, in-LLM tools)}

\begin{center}
\begin{tabular}{@{}lr@{}}
\toprule
\textbf{Per-call component} & \textbf{Cost} \\
\midrule
OpenAI: (280 + 8,000) input $\times$ \$0.75/M + 256 out $\times$ \$4.50/M + \$0.010 tool fee & \$0.01736 \\
Anthropic: (280 + 8,000) input $\times$ \$1/M + 256 out $\times$ \$5/M + 2$\times$\$0.010 tool fee & \$0.02956 \\
Perplexity Sonar: (280 input) \$1/M + 256 out \$1/M + \$0.005 request fee & \$0.00556 \\
Gemini (under quota): (280 + 3,500) \$0.30/M + 256 out \$2.50/M & \$0.00177 \\
Gemini (over quota): same + \$0.035 grounding fee & \$0.03677 \\
\bottomrule
\end{tabular}
\end{center}

\subsection{Per-run cost (44 prompts $\times$ 4 providers)}

\begin{center}
\begin{tabular}{@{}lrrr@{}}
\toprule
\textbf{Configuration} & \textbf{Conservative} & \textbf{Aggressive} \\
\midrule
OpenAI (176 $\div$ 4 = 44 calls)    & \$0.764 & \$1.452 \\
Anthropic (44 calls)                 & \$1.302 & \$1.302 \\
Perplexity Sonar (44 calls)          & \$0.244 & \$0.332 \\
Gemini under quota (44 calls)        & \$0.079 & --- \\
Gemini over quota (44 calls)         & ---      & \$1.620 \\
\midrule
\textbf{runCitationCheck total}      & \textbf{\$2.393} & \textbf{\$4.704} \\
\midrule
Ancillary calls (warm path)          & \$0.152 & \$0.152 \\
Ancillary calls (cold path)          & \$0.284 & \$0.284 \\
\midrule
\textbf{Per click, warm path}        & \textbf{\$2.55} & \textbf{\$4.86} \\
\textbf{Per click, cold path}        & \textbf{\$2.68} & \textbf{\$4.99} \\
\bottomrule
\end{tabular}
\end{center}

\subsection{Corrected margin table}

\begin{center}
\begin{tabular}{@{}lrrrrr@{}}
\toprule
\textbf{Customer} & \textbf{Pays} & \textbf{Cons. cost} & \textbf{Margin} & \textbf{Aggr. cost} & \textbf{Margin} \\
\midrule
Pro     & \$0.087 & \$2.55 & \costneg{-\$2.46} & \$4.86 & \costneg{-\$4.77} \\
Growth  & \$0.167 & \$2.55 & \costneg{-\$2.38} & \$4.86 & \costneg{-\$4.69} \\
Starter & \$0.346 & \$2.55 & \costneg{-\$2.20} & \$4.86 & \costneg{-\$4.51} \\
Pack    & \$0.500 & \$2.55 & \costneg{-\$2.05} & \$4.86 & \costneg{-\$4.36} \\
\bottomrule
\end{tabular}
\end{center}

\textbf{Headline:} every tier loses money on every click. Pro is the worst
because its credits are the cheapest. A Pro customer running Citation Check
5 times a month across 5 sites loses Flowblinq approximately \$87.50/month in
LLM spend on that single feature.

\section{Alternative architectures --- stripping web tools}

\begin{userturn}
Ok. If we don't engage web tools, how much does the citation check range
change by, in terms of costing?
\end{userturn}

\subsection{Mechanism}

Remove \texttt{tools: [\{ type: "web\_search", ... \}]} from the OpenAI,
Anthropic, and Gemini calls. Perplexity Sonar cannot have tools stripped ---
it is a search-native model with retrieval baked in. Two variants:

\begin{enumerate}[nosep]
  \item Keep Sonar with its built-in retrieval (4 providers, 176 calls)
  \item Drop Sonar entirely (3 providers, 132 calls)
\end{enumerate}

\subsection{Cost with tools off}

\begin{center}
\begin{tabular}{@{}lrrrr@{}}
\toprule
\textbf{Configuration} & \textbf{runCitationCheck} & \textbf{Ancillary} & \textbf{Per click} & \textbf{Reduction} \\
\midrule
Current (in-LLM tools, conservative) & \$2.39 & \$0.15 & \$2.55 & --- \\
Current (in-LLM tools, aggressive)   & \$4.70 & \$0.15 & \$4.86 & --- \\
Tools off, 4 providers (Sonar kept)  & \$0.41 & \$0.15 & \textbf{\$0.56} & \good{$-$78\% to $-$88\%} \\
Tools off, 3 providers (Sonar dropped) & \$0.16 & \$0.15 & \textbf{\$0.31} & \good{$-$88\% to $-$94\%} \\
\bottomrule
\end{tabular}
\end{center}

\subsection{Cost of stripping tools (non-cash)}

\begin{quote}
Stripping web tools changes what the citation check \textit{measures}. With
tools, it measures live discoverability --- what a real user would see when
the LLM actually searches the web right now. Without tools, it measures
parametric memory --- what the LLM remembers about the brand from its
training cutoff. The GEO product is positioned as ``improve your visibility
in AI search'' --- without tools, the citation check stops measuring
``AI search'' and starts measuring ``what AI remembers'', which is a
structural product change.
\end{quote}

\section{External search APIs as a replacement}

\begin{userturn}
Alright. Now, if instead of using the in llm tool, I use firecrawl /search,
then what happens? If I use Brave API, then what happens?
\end{userturn}

\subsection{Vendor pricing (verified live)}

\begin{center}
\begin{tabular}{@{}lllr@{}}
\toprule
\textbf{Vendor} & \textbf{Plan} & \textbf{Notes} & \textbf{\$/search} \\
\midrule
Brave Search    & Standard      & \$5/1k req; \$5/mo free credit $\approx$ 1,000 free queries; 50 QPS & \$0.00500 \\
Firecrawl \texttt{/search} & Hobby    & \$16/mo / 3,000 credits; 2 credits per search (no scrape) & \$0.01067 \\
Firecrawl \texttt{/search} & Standard & \$83/mo / 100,000 credits; 2 credits per search           & \$0.00166 \\
Firecrawl \texttt{/search} & Growth   & \$333/mo / 500,000 credits; 2 credits per search          & \$0.00133 \\
Firecrawl \texttt{/search} $+$scrape & Standard & 12 credits per search (full scraped content) & \$0.00996 \\
Firecrawl \texttt{/search} $+$scrape & Growth   & 12 credits per search                         & \$0.00800 \\
\bottomrule
\end{tabular}
\end{center}

\subsection{Architectural change}

Currently each of the 4 LLMs invokes its own web\_search tool independently
($\sim$132--176 total searches per run). The proposed architecture does
\textbf{one external search per prompt}, shared across all LLMs, and feeds
results into each LLM call as plain context. LLMs run with no tools.

Per run: \textbf{44 searches} (one per prompt), \textbf{176 LLM calls} (or
132 if Sonar is dropped, because Sonar's built-in retrieval makes it
redundant with external search).

\subsection{Total cost per click with external search}

\begin{center}
\begin{tabular}{@{}lrrrr@{}}
\toprule
\textbf{Configuration} & \textbf{Search} & \textbf{LLM} & \textbf{Anc.} & \textbf{Total} \\
\midrule
Current in-LLM tools (conservative)                                    & bundled & \$2.39 & \$0.15 & \$2.55 \\
Current in-LLM tools (aggressive)                                      & bundled & \$4.70 & \$0.15 & \$4.86 \\
Brave + 3 LLMs, snippets                                               & \$0.22  & \$0.43 & \$0.15 & \textbf{\$0.80} \\
Brave + 3 LLMs, free \$5 credit ($\le$22 runs/mo)                      & \$0.00  & \$0.43 & \$0.15 & \textbf{\$0.58} \\
Firecrawl Standard /search + 3 LLMs, snippets                          & \$0.07  & \$0.43 & \$0.15 & \textbf{\$0.66} \\
Firecrawl Growth /search + 3 LLMs, snippets                            & \$0.06  & \$0.43 & \$0.15 & \textbf{\$0.64} \\
Firecrawl Standard /search $+$scrape + 3 LLMs                          & \$0.44  & \$0.88 & \$0.15 & \textbf{\$1.47} \\
Firecrawl Growth /search $+$scrape + 3 LLMs                            & \$0.35  & \$0.88 & \$0.15 & \textbf{\$1.38} \\
\bottomrule
\end{tabular}
\end{center}

Cost reduction from current: 46\% (worst external) to 87\% (best external).

\section{The Profound/Athena claim retraction}

CoFounder initially wrote:

\begin{quote}
Profound and Athena both feed external search results too, so this is
industry standard, not a unique liability.
\end{quote}

\begin{userturn}
Profound and Athena both feed external search results too, so this is
industry standard, not a unique liability --- can you find me cited evidence
for this statement?
\end{userturn}

After searching, CoFounder retracted the claim. Actual findings:

\begin{itemize}[nosep]
  \item \textbf{Profound:} explicitly positions \textit{against} static API-based methods.
    Quote from their own blog (tryprofound.com): ``actively monitoring AI Search
    Engines like ChatGPT, Perplexity, and Copilot, rather than relying solely
    on static API model calls.'' This is the \textbf{opposite} of what was
    claimed. The strongest signal is that Profound uses direct interface
    monitoring, likely headless browser scraping of consumer interfaces.
  \item \textbf{Peec.ai} (a separate competitor) confirmed UI scraping: ``Peec AI
    uses UI scraping technology to simulate real user interactions with AI
    tools.''
  \item \textbf{AthenaHQ:} methodology is \textit{genuinely undisclosed}. No
    source verified whether they use LLM APIs, external search APIs, or
    browser automation.
\end{itemize}

\textbf{Corrected framing:} the Brave/Firecrawl approach being discussed is
a middle ground between ``trust the LLM API'' (current) and ``scrape the
consumer UI'' (likely Profound/Peec). Not an industry default. The liability
is larger than previously implied.

\section{Per-LLM search routing (and the death of the Bing API)}

\begin{userturn}
No. Lets go back to the previous thread --- using brave api for anthropic,
bing (if it has an api) for openAI and firecrawl api for Gemini.
\end{userturn}

\subsection{Critical finding}

Microsoft retired the Bing Search API on 2025-08-11. The replacement,
``Grounding with Bing Search'' via Azure AI Foundry, costs
\textbf{\$35 per 1,000 queries} --- a 2$\times$ to 5$\times$ increase over
the old S1--S5 tiers. This kills the OpenAI$\leftrightarrow$Bing pairing:

\begin{enumerate}[nosep]
  \item Bing's replacement is bundled, not standalone --- it requires an
    Azure AI Foundry agent wrapper
  \item At \$35/1k it is \textbf{3.5$\times$ more expensive} than OpenAI's
    own in-LLM web\_search tool (\$10/1k), so routing OpenAI through
    Bing-replacement is strictly worse than just keeping the in-LLM tool
\end{enumerate}

\subsection{Proposed routing cost (literal interpretation)}

\begin{center}
\begin{tabular}{@{}lrr@{}}
\toprule
\textbf{Leg} & \textbf{Per 44 prompts} \\
\midrule
Anthropic $\leftrightarrow$ Brave (\$0.005/search)             & \$0.220 + LLM \$0.201 = \$0.421 \\
OpenAI $\leftrightarrow$ Bing replacement (\$0.035/search)     & \$1.540 + LLM \$0.159 = \textbf{\$1.699} \\
Gemini $\leftrightarrow$ Firecrawl Growth (\$0.00133/search)   & \$0.059 + LLM \$0.071 = \$0.130 \\
\midrule
runCitationCheck total & \$2.250 \\
Ancillary              & \$0.152 \\
\textbf{Per click}     & \textbf{\$2.40} \\
\bottomrule
\end{tabular}
\end{center}

The Bing-replacement leg alone consumes 85\% of the search budget. Without
Bing this would be a $<$\$1/click configuration.

\subsection{Repair C --- the fix that works}

Keep the per-LLM routing principle, but drop Bing and reuse Brave across
OpenAI+Anthropic (with shared cache per prompt):

\begin{center}
\begin{tabular}{@{}lr@{}}
\toprule
\textbf{Component} & \textbf{Cost} \\
\midrule
Brave (1 search $\times$ 44, shared OpenAI+Anthropic) & \$0.220 \\
Gemini $\leftrightarrow$ Firecrawl (1 search $\times$ 44)   & \$0.059 \\
OpenAI LLM (+ Brave context)     & \$0.159 \\
Anthropic LLM (+ Brave context)  & \$0.201 \\
Gemini LLM (+ Firecrawl context) & \$0.071 \\
Ancillary & \$0.152 \\
\midrule
\textbf{Total per click} & \textbf{$\sim$\$0.862} \\
\bottomrule
\end{tabular}
\end{center}

\textbf{Note on Anthropic$\leftrightarrow$Brave:} Anthropic's own web\_search
server tool uses Brave Search as its underlying provider (publicly
documented). Calling Brave directly gives us the same content at
\$5/1k instead of Anthropic's \$10/1k markup, with no semantic loss.

\section{Firecrawl engine selection --- dead end}

\begin{userturn}
Can we get firecrawl to use specific search indices?
\end{userturn}

\textbf{No.} The Firecrawl \texttt{/search} endpoint parameter list
(docs.firecrawl.dev/api-reference/endpoint/search):

\begin{center}
\begin{tabular}{@{}ll@{}}
\toprule
\textbf{Parameter} & \textbf{Purpose} \\
\midrule
\texttt{query}            & search string \\
\texttt{limit}            & result count (default 5, max 100) \\
\texttt{sources}          & content vertical: web / news / images \\
\texttt{categories}       & content type: github / research / etc \\
\texttt{tbs}              & time filter \\
\texttt{location}         & geo locale \\
\texttt{country}          & ISO country code \\
\texttt{timeout}          & request timeout \\
\texttt{ignoreInvalidURLs}& filter dead URLs \\
\texttt{enterprise}       & ZDR options \\
\texttt{scrapeOptions}    & scraping config \\
\bottomrule
\end{tabular}
\end{center}

There is no \texttt{engine}, \texttt{provider}, \texttt{index}, or
\texttt{search\_type} parameter. Firecrawl does not publicly disclose its
underlying engine. The \texttt{sources} field selects content verticals
(web/news/images), not search engines.

The only vendors that expose engine selection as a first-class parameter are
SERP-aggregator APIs: SerpAPI, Serper, SearchAPI.io, ScaleSerp, Bright Data
SERP. For ``different engine per LLM'' the architecturally-correct vendor is
SerpAPI (\$0.015/search), not Firecrawl.

\section{MCP as a tool layer}

\begin{userturn}
Can these be served as tools within the LLM? [...] What if I move the
search to MCP? does that reduce costs?
\end{userturn}

\subsection{Tool calling support per provider}

\begin{center}
\begin{tabular}{@{}lllc@{}}
\toprule
\textbf{Provider} & \textbf{Mechanism} & \textbf{Status} & \textbf{Hard cap?} \\
\midrule
OpenAI (Responses API)    & \texttt{tools: [\{ type: "function", ... \}]}              & First-class & No \\
Anthropic (Messages)      & \texttt{tools: [\{ name, input\_schema \}]} + tool\_use    & First-class & \texttt{max\_uses} only on built-in web\_search, not custom tools \\
Gemini (generateContent)  & \texttt{tools: [\{ functionDeclarations: [...] \}]}       & First-class & No \\
Perplexity Sonar          & Limited (\texttt{sonar-pro} only; base \texttt{sonar} no) & Not supported on our model & --- \\
\bottomrule
\end{tabular}
\end{center}

\subsection{Multi-turn cost problem}

Every tool invocation requires at least \textbf{two LLM round-trips}: the
first to receive the tool\_use block, the second to send the tool\_result
and receive the final answer. Each round-trip re-sends the full context.
Without a hard cap, the model may call the tool 2--5 times. Cost and
variance both return.

\begin{center}
\begin{tabular}{@{}lrr@{}}
\toprule
\textbf{Strategy} & \textbf{Round-trips} & \textbf{Per-prompt cost (Haiku)} \\
\midrule
Current in-LLM tool (capped at 2 searches) & 1 & \$0.0296 \\
Stuff results in prompt                    & 1 & \$0.0096 \\
Custom tool (1 search avg)                 & 2 & \$0.0102 \\
Custom tool (2 searches avg)               & 3 & \$0.0153 \\
Custom tool (uncapped, 4-deep)             & 5 & \$0.030+ \\
\bottomrule
\end{tabular}
\end{center}

\subsection{MCP cost analysis}

\textbf{Short answer: MCP does not reduce LLM token costs.} It is a protocol
standardization, not a pricing tier.

\begin{center}
\begin{tabular}{@{}lcc@{}}
\toprule
\textbf{Cost component} & \textbf{Without MCP} & \textbf{With MCP} \\
\midrule
Token billing for input/output     & Same & \textbf{Same} \\
Tool definition overhead (~80 tok) & Yes  & Yes \\
Multi-turn round-trips             & Yes  & Yes (internally, even in Claude's remote MCP) \\
Search vendor cost                 & Direct (we pay Brave) & Direct (we pay Brave through the wrapper) \\
LLM provider fee for MCP vs custom tool & n/a & No difference \\
\bottomrule
\end{tabular}
\end{center}

What MCP \textit{does} provide (architectural, not cost):

\begin{itemize}[nosep]
  \item Write one search tool, reuse across all MCP-compatible LLMs
  \item Server-side caching shared across providers in one runtime
  \item Swap search backends without touching LLM code
  \item Clean separation between the search layer and the LLM layer
\end{itemize}

The one angle where MCP \textit{could} reduce cost is cross-customer search
deduplication via a shared MCP server cache. But the same savings are
available by building an application-level shared search cache (Upstash
Redis keyed on \texttt{(query, day)}), which does not require MCP.

\textbf{Additional blocker:} only OpenAI and Anthropic have MCP client
support in their server-side APIs. Gemini has MCP support only in Gemini
CLI and Gemini Code Assist, not in \texttt{generateContent}. Perplexity
Sonar has no MCP client support. So even with an MCP wrapper, we would
still need per-provider glue for half the providers.

\subsection{The one consumer-side scenario (and why it doesn't apply)}

If citation checks were run interactively by users inside Claude Desktop
with their personal \$20/mo Claude Pro subscription, MCP could let those
users invoke our search tool without us paying any LLM tokens. But:

\begin{itemize}[nosep]
  \item Citation check is an automated SaaS feature, not an interactive conversation
  \item It cannot run on a cron against a desktop app
  \item It only covers Claude, not the other three providers
\end{itemize}

So MCP's consumer-subscription bypass doesn't apply to our product.

\section{Prompt caching}

\subsection{Per-provider specifications (verified live)}

\begin{center}
\begin{tabular}{@{}lllll@{}}
\toprule
 & \textbf{OpenAI gpt-5.4-mini} & \textbf{Anthropic Haiku 4.5} & \textbf{Gemini 2.5 Flash} \\
\midrule
Mode & Automatic, no opt-in & Explicit (\texttt{cache\_control}) & Implicit + explicit \\
Minimum prefix & 1,024 tok & \warn{\textbf{4,096 tok}} & 1,024 tok (explicit) \\
Cache write cost & base rate (no premium) & \$1.25/M (5m) or \$2/M (1h) & base rate \\
Cache read cost & \$0.075/M (\good{-90\%}) & \$0.10/M (\good{-90\%}) & $\sim$\$0.075/M ($\sim$-75\%) \\
TTL & 5--10 min, \texttt{prompt\_cache\_retention} & 5 min default, 1h for 2$\times$ write & 1 hour default \\
Breakpoints per req & n/a (auto) & up to 4 & 1 \\
\bottomrule
\end{tabular}
\end{center}

\subsection{Why citation-checker.ts gets zero caching benefit as-is}

The stable prefix across all 176 calls in one run is the
\texttt{CITATION\_SYSTEM\_PROMPT} at $\sim$250 tokens. This is
\textbf{well under} every provider's minimum:

\begin{center}
\begin{tabular}{@{}llll@{}}
\toprule
\textbf{Provider} & \textbf{Min cacheable} & \textbf{Our prefix} & \textbf{Eligible?} \\
\midrule
OpenAI              & 1,024 & $\sim$250 & No \\
Anthropic Haiku 4.5 & \warn{\textbf{4,096}} & $\sim$250 & No \\
Gemini (explicit)   & 1,024 & $\sim$250 & No \\
Gemini (implicit)   & --    & $\sim$250 & Maybe, no guarantee \\
\bottomrule
\end{tabular}
\end{center}

\textbf{Pure citation-checker.ts receives essentially zero benefit from
prompt caching.} This is a correction on an earlier overstatement.

\subsection{Path to actually unlock caching}

Strategy C: pad the stable prefix to 4,096 tokens with load-bearing content
(output format guides, sentiment rubric, per-site brand context). Mark
\texttt{cache\_control: \{ type: "ephemeral" \}}.

\begin{center}
\begin{tabular}{@{}lrrr@{}}
\toprule
\textbf{Provider} & \textbf{Without caching (176 $\times$ 4K)} & \textbf{With caching (1 write + 175 reads)} & \textbf{Savings} \\
\midrule
OpenAI             & \$0.541 & \$0.057 & \good{\$0.484 (89\%)} \\
Anthropic Haiku    & \$0.721 & \$0.076 & \good{\$0.645 (89\%)} \\
Gemini Flash       & \$0.216 & \$0.026 & \good{\$0.190 (88\%)} \\
\midrule
\textbf{Combined}  & \$1.479 & \textbf{\$0.160} & \good{\textbf{\$1.319/run}} \\
\bottomrule
\end{tabular}
\end{center}

This requires the padding to be functionally useful, not noise.

\section{Preemptive tool injection}

\subsection{Mechanism}

Run the search \textit{outside} the LLM, then construct a synthetic prior
tool exchange so the model sees ``I already called web\_search, here's the
tool\_result, now answer''. Single round-trip.

Anthropic code shape:

\begin{lstlisting}[language=JavaScript]
const results = await brave.search({ query: prompt, count: 10 });

await client.messages.create({
  model: "claude-haiku-4-5-20251001",
  max_tokens: 256,
  system: CITATION_SYSTEM_PROMPT,
  tools: [{
    name: "web_search",
    description: "Search the web for current information",
    input_schema: { type: "object", properties: { query: { type: "string" } } }
  }],
  messages: [
    { role: "user", content: prompt },
    { role: "assistant", content: [
      { type: "tool_use", id: "toolu_pre_001",
        name: "web_search", input: { query: prompt } }
    ]},
    { role: "user", content: [
      { type: "tool_result",
        tool_use_id: "toolu_pre_001",
        content: formatResults(results) }
    ]}
  ]
});
\end{lstlisting}

\subsection{Comparison matrix}

\begin{center}
\begin{tabular}{@{}lccc@{}}
\toprule
\textbf{Property} & \textbf{Stuff results} & \textbf{Full tool calling} & \textbf{Preemptive injection} \\
\midrule
Round-trips                                    & 1           & 2+           & \textbf{1} \\
LLM-formulated query                           & No          & Yes          & No (we pick) \\
Model treats evidence as search output         & No          & Yes          & \textbf{Yes} \\
Model can chase a follow-up search             & No          & Yes          & \textbf{Yes (rare)} \\
Predictable cost                               & Yes         & No           & \textbf{Yes} \\
Token overhead vs stuff-results                & 0           & +500--2000   & \textbf{$\sim$+130} \\
\bottomrule
\end{tabular}
\end{center}

\subsection{Cost of preemptive injection with shared Brave across 3 providers}

\begin{center}
\begin{tabular}{@{}lrr@{}}
\toprule
\textbf{Component} & \textbf{Per call} & \textbf{$\times$ 44} \\
\midrule
Brave search (shared across OpenAI+Anthropic+Gemini) & \$0.005 & \$0.220 \\
OpenAI gpt-5.4-mini (3,390 in, 256 out)              & \$0.00370 & \$0.163 \\
Anthropic Haiku 4.5 (3,390 in, 256 out)              & \$0.00467 & \$0.205 \\
Gemini 2.5 Flash (3,390 in, 256 out)                 & \$0.00166 & \$0.073 \\
\midrule
runCitationCheck & & \$0.661 \\
Ancillary         & & \$0.152 \\
\textbf{Per click} & & \textbf{$\sim$\$0.81} \\
\bottomrule
\end{tabular}
\end{center}

\section{50-credit bump analysis}

\begin{userturn}
I want to actually increase the cost to 50 credits per citation check\ldots
where does that put us?
\end{userturn}

\subsection{Customer-side math}

A 10$\times$ price increase. Effective price per click:

\begin{center}
\begin{tabular}{@{}lrr@{}}
\toprule
\textbf{Credit source} & \textbf{At 5 credits} & \textbf{At 50 credits} \\
\midrule
Pro          & \$0.087 & \$0.864 \\
Growth       & \$0.167 & \$1.671 \\
Starter      & \$0.346 & \$3.462 \\
One-off pack & \$0.500 & \$5.000 \\
\bottomrule
\end{tabular}
\end{center}

\subsection{Where 50 credits breaks customer budgets}

At 50 credits per check, citation check consumes 80--95\% of every
customer's monthly credit budget. Growth and Pro can no longer support
daily monitoring on every site they are nominally allowed to track:

\begin{center}
\begin{tabular}{@{}lrrrrr@{}}
\toprule
\textbf{Tier} & \textbf{Sites} & \textbf{Cadence} & \textbf{Wanted checks/mo} & \textbf{Affordable at 50cr} & \textbf{Gap} \\
\midrule
Starter (weekly) & 5  & 4/mo/site     & 20          & 28  & +8 spare \\
Growth (daily)   & 10 & 30/mo/site    & 300         & 149 & \warn{$-$151 short} \\
Pro (daily)      & 20 & 30/mo/site    & 600         & 578 & \warn{$-$22 short} \\
\bottomrule
\end{tabular}
\end{center}

\section{Aditya's 2/day usage model}

\begin{userturn}
So the way I think about this is is that the user will want to do a citation
check 2 times a day, this means 20 citation checks per day for growth --- 1000
credits/day --- in 25 days --- 25000 credits\ldots and as many audits/week as
the user wants.
\end{userturn}

\subsection{What the model implies per tier}

\begin{center}
\begin{tabular}{@{}lrrrrrr@{}}
\toprule
\textbf{Tier} & \textbf{Sites} & \textbf{Checks/day} & \textbf{Checks/mo} & \textbf{Credits needed @ 50cr} & \textbf{Current} & \textbf{Shortfall} \\
\midrule
Starter & 5  & 10 & 250   & 12,500 & 1,430  & \warn{$-$11,070 (8.7$\times$)} \\
Growth  & 10 & 20 & 500   & 25,000 & 7,465  & \warn{$-$17,535 (3.35$\times$)} \\
Pro     & 20 & 40 & 1,000 & 50,000 & 28,915 & \warn{$-$21,085 (1.73$\times$)} \\
\bottomrule
\end{tabular}
\end{center}

Every paid tier's credit allocation breaks. Tier credit budgets do not
sustain the intended usage pattern at 50 credits per check.

\subsection{2$\times$-cost-on-Pro formula}

\begin{center}
\begin{tabular}{@{}lrrr@{}}
\toprule
\textbf{Infrastructure path} & \textbf{Cost/click} & \textbf{2$\times$ cost target} & \textbf{Credits (@ \$0.01726/cr)} \\
\midrule
Current in-LLM tools (aggressive)             & \$4.86 & \$9.72 & 563 \\
Current in-LLM tools (conservative)           & \$2.55 & \$5.10 & 295 \\
Tools off, 3 providers (Sonar dropped)        & \$0.31 & \$0.62 & 36 \\
Stuff Firecrawl Growth snippets               & \$0.64 & \$1.28 & 74 \\
Stuff Brave snippets                          & \$0.80 & \$1.60 & 93 \\
Preemptive injection $+$ Brave (warm)         & \$0.86 & \$1.72 & $\sim$100 \\
Preemptive injection $+$ caching $+$ Brave    & \$0.86 & \$1.72 & $\sim$100 \\
Full stack ($+$ response cache, 2/day pattern) & \$0.46 & \$0.92 & $\sim$53 \\
Full stack average (mixed usage)              & \$0.25 & \$0.50 & $\sim$29 \\
\bottomrule
\end{tabular}
\end{center}

\section{Sonar retention and MCP reassessment}

\begin{userturn}
I don't want to drop sonar\ldots Perplexity has enough penetration in the
markets\ldots also, just a second thought? What if I move the search to MCP?
does that reduce costs?
\end{userturn}

\subsection{Sonar kept --- updated 4-provider cost}

\begin{center}
\begin{tabular}{@{}lr@{}}
\toprule
\textbf{Component} & \textbf{Cost} \\
\midrule
Brave search $\times$ 44 (shared OpenAI+Anthropic+Gemini) & \$0.220 \\
OpenAI preemptive injection                                & \$0.163 \\
Anthropic preemptive injection                             & \$0.206 \\
Gemini preemptive injection                                & \$0.073 \\
\textbf{Sonar (native retrieval, 44 calls)}                & \textbf{\$0.246} \\
Ancillary (warm)                                           & \$0.152 \\
\midrule
\textbf{Per click (warm, no cache)}                        & \textbf{$\sim$\$1.06} \\
\textbf{Per click (prompt caching)}                        & \textbf{$\sim$\$0.96} \\
\textbf{Per click (2/day blended with response caching)}   & \textbf{$\sim$\$0.58} \\
\bottomrule
\end{tabular}
\end{center}

Sonar adds $\sim$\$0.25 per click. That is the price of measuring
Perplexity's retrieval behavior.

\subsection{2$\times$-cost-on-Pro updated for Sonar inclusion}

\begin{center}
\begin{tabular}{@{}lrrr@{}}
\toprule
\textbf{Path} & \textbf{Cost/click} & \textbf{2$\times$ target} & \textbf{Credits @ \$0.01726/cr} \\
\midrule
Warm (no response cache)            & \$1.06 & \$2.12 & $\sim$123 \\
Blended on 2/day (cache hot)        & \$0.58 & \$1.16 & $\sim$67 \\
Average across mixed usage          & \$0.35 & \$0.70 & $\sim$41 \\
\bottomrule
\end{tabular}
\end{center}

\subsection{MCP verdict --- no cost reduction}

CoFounder verified that MCP does not reduce LLM token costs. Token billing
is identical whether a tool is exposed via MCP, custom function calling, or
preemptive injection. MCP's real value is architectural: single tool
implementation, server-side caching, clean separation. But Gemini does not
support MCP clients in its server-side API, and Sonar does not support MCP
at all, so even with MCP we still need per-provider glue for half the
providers.

\textbf{Decision: skip MCP for the cost optimization pass. Consider it later
as a distribution/integration lever.}

\section{Usage-based pricing model}

\begin{userturn}
I really like this idea though: ``E. Usage-based pricing: fixed low sub $+$
metered citation check. \$99 base $+$ \$0.50/check for everyone.''
\end{userturn}

CoFounder flagged the earlier ``Profound/Athena converging on this'' claim
as overstated (not verified), and grounded the recommendation on the
observable base-plus-metered pattern in Vercel AI Gateway, OpenAI API,
Replicate, Anthropic API, Modal, Together, and Fireworks --- i.e., the
pattern is well-established in general AI SaaS even if specific competitor
adoption is undetermined.

\subsection{Why usage-based fits this product}

\begin{enumerate}[nosep]
  \item Cost is dominated by variable unit cost ($\sim$80\% LLM + search API fees scaling 1:1 with clicks)
  \item Usage varies wildly by customer type; tier caps fail for both ends
  \item Citation check is a clean discrete action, easy to meter
\end{enumerate}

\section{Final pricing decision}

\begin{userturn}
First. The subscription pricing stays at --- \$99 starter, 249\$ growth,
499\$ pro. Lets call them exactly what they are. Second --- verbatim and in
full fidelity please draft a detailed latex document of this conversation in
geo/docs/ folder. Third, the pricing solution is to look at our cost,
multiply it by 1.5x and then arrive at dollar value pricing for overage. We
adjust our costing there. Fourth, we should stage the architecture in a
simple script and verify that it works in the case of one LLM. this should
be a standalone script. Please write the TS for that script.
\end{userturn}

\subsection{Decision 1: Fixed subscription prices stay}

\textbf{Starter \$99 / Growth \$249 / Pro \$499 per month.} Names unchanged.
No tier rename, no price shift.

\subsection{Decision 2: Overage pricing formula}

\textbf{Overage rate $=$ variable cost per check $\times$ 1.5.}

Applied to every cost architecture:

\begin{center}
\begin{tabular}{@{}lrrr@{}}
\toprule
\textbf{Architecture} & \textbf{Cost/click} & \textbf{$\times$ 1.5} & \textbf{Rounded} \\
\midrule
Current (in-LLM tools, aggressive)                               & \$4.86 & \$7.29 & \$7.50 \\
Current (in-LLM tools, conservative)                             & \$2.55 & \$3.83 & \$3.75 \\
Preemptive injection $+$ Brave $+$ Sonar (warm, no cache)        & \$1.06 & \$1.59 & \textbf{\$1.60} \\
Preemptive injection $+$ Brave $+$ Sonar $+$ prompt caching      & \$0.96 & \$1.44 & \$1.45 \\
Preemptive injection $+$ full response cache $+$ Sonar           & \$0.58 & \$0.87 & \$0.90 \\
Full stack across mixed usage                                    & \$0.35 & \$0.53 & \$0.55 \\
\bottomrule
\end{tabular}
\end{center}

\textbf{Lock-in at ship time: overage \$1.60/check} (based on the warm,
no-response-cache path). Recalibrate downward once response caching ships
and real-world cost is measured from provider invoices.

\subsection{Decision 3: Free allowances per tier}

Working backward from each base sub budget, assuming \$0.40/audit and
\$1.06/check variable cost, with a platform margin buffer:

\begin{center}
\begin{tabular}{@{}lrrrrr@{}}
\toprule
\textbf{Tier} & \textbf{Base} & \textbf{Platform margin buffer} & \textbf{Audit budget} & \textbf{Check budget} & \textbf{Free checks} \\
\midrule
Starter (5 sites, weekly) & \$99  & \$30  & \$8 (20 audits) & \$61  & \textbf{50} \\
Growth  (10 sites, daily) & \$249 & \$60  & \$100 (250)     & \$89  & \textbf{80} \\
Pro     (20 sites, daily) & \$499 & \$100 & \$200 (500)     & \$199 & \textbf{200} \\
\bottomrule
\end{tabular}
\end{center}

\subsection{Final pricing table}

\begin{center}
\begin{tabular}{@{}llrrr@{}}
\toprule
\textbf{Plan} & \textbf{Base/mo} & \textbf{Sites} & \textbf{Free checks} & \textbf{Overage} \\
\midrule
Free    & \$0   & 1  & 1 lifetime & no overage (must upgrade) \\
\textbf{Starter} & \textbf{\$99}  & 5  & \textbf{50}  & \textbf{\$1.60/check} \\
\textbf{Growth}  & \textbf{\$249} & 10 & \textbf{80}  & \textbf{\$1.60/check} \\
\textbf{Pro}     & \textbf{\$499} & 20 & \textbf{200} & \textbf{\$1.60/check} \\
\bottomrule
\end{tabular}
\end{center}

\subsection{Margin verification at 2/day usage model}

\begin{center}
\begin{tabular}{@{}lrrrrrr@{}}
\toprule
\textbf{Tier} & \textbf{Checks/mo} & \textbf{Free} & \textbf{Overage} & \textbf{Revenue} & \textbf{Variable cost} & \textbf{Margin} \\
\midrule
Starter (5 $\times$ 1/day)  & 125   & 50  & \$120  & \$219  & \$140.50 & \good{+\$78.50} \\
Growth (10 $\times$ 2/day)  & 500   & 80  & \$672  & \$921  & \$630    & \good{+\$291} \\
Pro (20 $\times$ 2/day)     & 1,000 & 200 & \$1,280 & \$1,779 & \$1,260  & \good{+\$519} \\
\bottomrule
\end{tabular}
\end{center}

Every tier is profitable at the designed 1.5$\times$ cost markup. Heavy
users generate disproportionate absolute margin.

\subsection{Decision 4: Architecture staging script}

A standalone TypeScript verification script has been written to
\texttt{geo/scripts/verify-preemptive-injection.ts}. The script:

\begin{itemize}[nosep]
  \item Takes a citation prompt as CLI argument
  \item Runs one Brave Search call externally
  \item Constructs a synthetic assistant tool\_call $+$ tool result exchange
  \item Sends one round-trip call to OpenAI \texttt{gpt-5.4-mini}
  \item Measures latency, token usage, and dollar cost breakdown
  \item Prints a per-stage report and extrapolates to a full 44-prompt run
\end{itemize}

Usage:

\begin{lstlisting}[language=bash]
BRAVE_API_KEY=... OPENAI_API_KEY=... \
  npx tsx geo/scripts/verify-preemptive-injection.ts \
  "What are the best dental clinics in Mumbai?"
\end{lstlisting}

The script has no imports from \texttt{lib/}, does not touch the database,
does not consult \texttt{config.ts}, and does not charge credits. It is a
clean-room verification of the architecture on a single LLM before wiring
it into production.

\section{Next steps and outstanding items}

\subsection{Hard dependencies before pricing ships}

\begin{enumerate}[nosep]
  \item Run \texttt{verify-preemptive-injection.ts} against live Brave and
    OpenAI credentials; confirm measured cost is within $\pm$20\% of analytical estimate
  \item Extend the script to Anthropic and Gemini in follow-up verifications
  \item Audit the GEO audit pipeline (\texttt{/api/sites/[id]/regenerate}) for
    its own cost per run --- current \$0.40/audit estimate is unverified
  \item Fix \texttt{rephraseSeeds} \texttt{temperature: 0.7} in
    \texttt{citation-prompt-generator.ts:1222} for prompt determinism if
    response caching is to work
\end{enumerate}

\subsection{Spec follow-ups}

\begin{itemize}[nosep]
  \item \textbf{TS-077 Architecture:} preemptive injection + Brave + Sonar retention + prompt determinism fix
  \item \textbf{TS-078 Response caching:} Upstash Redis 24h cache on \texttt{(provider, model, prompt\_hash)}
  \item \textbf{TS-079 Pricing migration:} Stripe metered billing pipeline + tier migration + overage guardrails
  \item \textbf{TS-080 Audit cost audit:} same treatment for the GEO audit pipeline
\end{itemize}

\subsection{Open questions}

\begin{itemize}[nosep]
  \item Actual response cache hit rate in production once prompt determinism is fixed
  \item Real-world token usage for preemptive injection (verification script will measure)
  \item Audit cost verification (currently assumed \$0.40/audit)
  \item Whether the 4,096-token padded system prompt for Anthropic Haiku caching is load-bearing or noise
  \item Migration strategy for existing credit-model customers
\end{itemize}

\section{Acknowledgments and corrections}

CoFounder made two factual errors during this session that Aditya caught
and corrected:

\begin{enumerate}[nosep]
  \item \textbf{The Profound/Athena ``industry standard'' claim} was
    unverified and later retracted after searching. Profound explicitly
    positions against API-based approaches and likely uses headless browser
    scraping of consumer interfaces. AthenaHQ's methodology is undisclosed.
    Peec.ai is confirmed to use UI scraping. The Brave/Firecrawl approach
    is \textit{not} an industry default --- it is a middle ground.
  \item \textbf{The initial ``\$1.20 per click'' customer revenue figure}
    used \texttt{BULK\_CREDIT\_PRICE\_INR = 20} (\$0.24/credit), which only
    applies to the bulk CSV audit path. The citation check charge draws
    from the fungible credit pool at subscription-tier prices ranging from
    \$0.0173/credit (Pro) to \$0.10/credit (pack). Actual customer revenue
    per 5-credit click: \$0.087 to \$0.500. The correction made the margin
    gap wider than originally reported.
\end{enumerate}

Both corrections improved the analysis. They are recorded here to prevent
their re-introduction in follow-up specifications.

\vspace{1em}

\begin{center}
\rule{0.5\textwidth}{0.4pt}
\end{center}

\vspace{0.5em}

\begin{center}\small
\textit{End of document. Associated artifacts:}\\
\texttt{geo/scripts/verify-preemptive-injection.ts} (verification script)\\
\texttt{geo/lib/services/citation-checker.ts} (audited source)\\
\texttt{geo/app/api/sites/[id]/citation-check/route.ts} (route wrapper)\\
\texttt{geo/lib/config.ts} (pricing and tier configuration)
\end{center}

\end{document}
